\section{Projektvorgehen}
\label{sec:Projektvorgehen}

\subsection{Planungsphase}
\label{sec:Planungsphase}
Die Umsetzung des Projektes umfasst insgesamt 80 Arbeitsstunden. Diese wurden im Vorfeld in verschiedene Phasen aufgeteilt, die während der Projektentwicklung durchlaufen werden. Eine grobe Übersicht der Zeitplanung sowie die Hauptphasen können der \TabelleRef{tab:GroberZeitplan} entnommen werden. % Frage: Zeitform - werden oder wurden?
Diese Hauptphasen lassen sich in kleinere Arbeitsschritte unterteilen, um den Fortschritt und die Aufgabenverteilung genauer darzustellen. Eine detaillierte Übersicht der einzelnen Phasen wird im \Anhang{app:Zeitplanung} aufgeschlüsselt.

\tabelle{Grober Zeitplan}{tab:GroberZeitplan}{GroberZeitplan}

\subsection{Ressourcenplanung}
\label{sec:Ressourcenplanung}
Die Auswahl der verwendeten Software erfolgte unter der Maßgabe, auf bereits lizenzierte oder frei verfügbare (Open Source) Software zurückzugreifen, um die Projektkosten zu minimieren. Eine detaillierte Auflistung der eingesetzten Tools und Ressourcen befindet sich im \Anhang{app:Ressourcen}.

\subsection{Entwicklungsprozess}
\label{sec:Entwicklungsprozess}
Im Rahmen dieses Projektes wurde das erweiterte Wasserfallmodell gewählt. Dieses bietet eine klar strukturierte und sequenzielle Vorgehensweise, die den Anforderungen und der Komplexität dieses Projektes optimal entspricht. 
Das erweiterte Wasserfallmodell basiert auf dem klassischen Wasserfall-modell, welches durch seine Phasenstruktur – \textbf{Anforderungsanalyse}, \textbf{Design}, \textbf{Implementierung}, \textbf{Test} und \textbf{Wartung} – einen linearen Projektablauf gewährleistet. Darüber hinaus bietet das erweiterte Modell die Flexibilität, frühere Phasen zu revidieren und Anpassungen vorzunehmen, falls während der Entwicklung unerwartete Änderungen oder Erkenntnisse auftreten. Diese erweiterte Version erlaubt somit gezielte Rücksprünge in frühere Phasen, ohne die gesamte Prozesskette zu unterbrechen. Dies verringert die Risiken, die durch Fehlspezifikationen oder unvorhergesehene Anforderungen entstehen könnten.
% Quelle notwendig?